% ============================================================================
%  COTALIA: sistema de diseño editorial
%  Tipografía: Inter (texto y titulares) + IBM Plex Mono (etiquetas, datos)
%  Paleta: tinta (verde casi negro), papel cálido, acento laterita (cobre)
% ============================================================================
\usepackage{fontspec}
\usepackage[spanish, es-noshorthands, es-tabla, es-nodecimaldot]{babel}
\usepackage[a4paper, top=27mm, bottom=26mm, inner=24mm, outer=24mm, headheight=12mm, headsep=9mm, footskip=12mm]{geometry}
\usepackage{xcolor}
\usepackage{graphicx}
\usepackage{tikz}
\usetikzlibrary{positioning, arrows.meta, calc, backgrounds}
\usepackage{pgfgantt}
\usepackage{booktabs, tabularx, array, multirow, colortbl, longtable}
\usepackage{catchfile}
\usepackage{enumitem}
\usepackage{caption}
\usepackage{float}
\usepackage{fancyhdr}
\usepackage[nobottomtitles*]{titlesec}
\usepackage{titletoc}
\usepackage{ragged2e}
\usepackage{lastpage}
\usepackage{qrcode}
\usepackage{microtype}
\usepackage{needspace}
\usepackage{etoolbox}
\usepackage[outline]{contour}
\contourlength{1.2pt}
\usepackage[hidelinks]{hyperref}
\def\UrlBreaks{\do\/}

% ---------------------------------------------------------------- paleta
\definecolor{tinta}{HTML}{0C1A17}       % casi negro, subtono verde
\definecolor{grafito}{HTML}{2E3834}     % texto
\definecolor{pizarra}{HTML}{717B77}     % texto secundario
\definecolor{linea}{HTML}{D8DBD6}       % filetes
\definecolor{niebla}{HTML}{F3F2EE}      % paneles
\definecolor{cobre}{HTML}{C4622D}       % acento: laterita
\definecolor{cobreclaro}{HTML}{E9A67F}
\definecolor{dosel}{HTML}{2F6B4F}       % verde secundario

% ---------------------------------------------------------------- tipografía
\setmainfont{Inter-Regular.otf}[BoldFont=Inter-SemiBold.otf, ItalicFont=Inter-Italic.otf,
  BoldItalicFont=Inter-SemiBoldItalic.otf, Scale=0.94, RawFeature=-calt]
\setmonofont{IBMPlexMono-Regular.otf}[BoldFont=IBMPlexMono-Medium.otf, Scale=0.94]
\newfontfamily\thin{Inter-Thin.otf}[RawFeature=-calt]
\newfontfamily\xlight{Inter-ExtraLight.otf}[RawFeature=-calt]
\newfontfamily\light{Inter-Light.otf}[RawFeature=-calt]
\newfontfamily\medium{Inter-Medium.otf}[RawFeature=-calt]
\newfontfamily\semi{Inter-SemiBold.otf}[RawFeature=-calt]
\newfontfamily\monomed{IBMPlexMono-Medium.otf}
\newfontfamily\monoreg{IBMPlexMono-Regular.otf}
\newcommand{\tnum}{\addfontfeature{Numbers=Monospaced}}

% ---------------------------------------------------------------- matemáticas (Fira Math, sans, a juego con Inter)
\usepackage{mathtools}
\usepackage{unicode-math}
\setmathfont{FiraMath-Regular.otf}[Scale=0.94]
% número con coma decimal dentro de fórmulas: 0,25 -> 0{,}25 (sin espacio de puntuación tras la coma)
\ExplSyntaxOn
\NewDocumentCommand{\mnum}{m}{ \tl_set:Ne \l_tmpa_tl {#1} \tl_replace_all:Nnn \l_tmpa_tl {,} {{,}} \tl_use:N \l_tmpa_tl }
\ExplSyntaxOff
% fórmulas tras una línea corta: mismo espacio que tras una línea larga
\makeatletter
\g@addto@macro\normalsize{\setlength\abovedisplayshortskip{\abovedisplayskip}\setlength\belowdisplayshortskip{\belowdisplayskip}}
\makeatother
\newtagform{cotalia}{\monoreg\fontsize{7.6}{9}\selectfont\color{pizarra}(}{)}
\usetagform{cotalia}
\setlength{\jot}{6pt}

% ---------------------------------------------------------------- logotipos de la tecnología usada
\linespread{1.32}
\setlength{\parskip}{0.55em}
\setlength{\parindent}{0pt}
\AtBeginDocument{\color{grafito}}
\frenchspacing
\AtBeginDocument{\RaggedRight}
\setlength{\emergencystretch}{2.5em}
\renewcommand{\bottomtitlespace}{0.18\textheight}

% etiqueta superior (eyebrow): mono, mayúsculas, espaciada
\newcommand{\eyebrow}[2][cobre]{\leavevmode{\hyphenpenalty=10000\monomed\fontsize{6.8}{9}\selectfont\addfontfeature{LetterSpace=14}\color{#1}\MakeUppercase{#2}}}
\newcommand{\lede}[1]{\leavevmode{\light\fontsize{13.2}{19.8}\linespread{1}\selectfont\raggedright\hyphenpenalty=10000\exhyphenpenalty=10000\color{tinta}#1\par}\vspace{0.7em}}
\newcommand{\datos}[1]{{\monoreg\small#1}}

% ---------------------------------------------------------------- numeración 01, 02...
\makeatletter
\renewcommand{\thesection}{\two@digits{\value{section}}}
\renewcommand{\thesubsection}{\thesection.\arabic{subsection}}
\renewcommand{\thefigure}{\two@digits{\value{figure}}}
\renewcommand{\thetable}{\two@digits{\value{table}}}
\renewcommand{\theequation}{\two@digits{\value{equation}}}
\newcommand{\numpag}{\two@digits{\value{page}}}
\newcommand{\indice}{\@starttoc{toc}}
\makeatother

% ---------------------------------------------------------------- títulos
\newcommand{\sectionbreak}{\clearpage}
\titleformat{\section}[display]
  {\normalfont}
  {\thin\fontsize{54}{54}\selectfont\color{cobre}\thesection}
  {6pt}
  {\xlight\fontsize{29}{34}\linespread{1}\selectfont\color{tinta}}
  [\vspace{2pt}{\color{linea}\rule{\linewidth}{0.4pt}}]
\titlespacing*{\section}{0pt}{-4pt}{12pt}
\titleformat{\subsection}[display]
  {\semi\fontsize{12.5}{16}\selectfont\color{tinta}}
  {\eyebrow{\thesubsection}}
  {3pt}
  {}
\titlespacing*{\subsection}{0pt}{16pt}{5pt}

% índice
\contentsmargin{0pt}
\makeatletter
\titlecontents{section}[2.6em]{\addvspace{7pt}}
  {\contentslabel[{\monomed\fontsize{8}{10}\selectfont\color{cobre}\thecontentslabel}]{2.6em}\medium\color{tinta}}
  {\hspace*{-2.6em}\medium\color{tinta}}
  {\hfill{\monoreg\small\color{pizarra}\two@digits{\thecontentspage}}}
  [\vspace{3pt}{\color{linea}\hrule height 0.3pt}]
\makeatother

% ---------------------------------------------------------------- logo (vectorial)
% Marca: tres curvas de nivel excéntricas y abiertas forman una "C"; el punto de cobre es la cota.
\newcommand{\CotaliaMarca}[3]{% #1 escala, #2 color del trazo, #3 color del punto
  \begin{tikzpicture}[scale=#1, line cap=round, baseline={(0,-0.32)}]
    \draw[#2, line width=#1*2.15pt] (0,0) ++(42:1) arc[start angle=42, end angle=318, radius=1];
    \draw[#2, line width=#1*2.15pt] (0.10,0.05) ++(48:0.70) arc[start angle=48, end angle=312, radius=0.70];
    \draw[#2, line width=#1*2.15pt] (0.19,0.10) ++(58:0.40) arc[start angle=58, end angle=302, radius=0.40];
    \fill[#3] (0.27,0.14) circle (0.135);
  \end{tikzpicture}}
\newcommand{\CotaliaLogo}[3]{% #1 escala, #2 color, #3 color del punto
  \mbox{\CotaliaMarca{#1}{#2}{#3}\hspace{\fpeval{#1*0.42}cm}%
  {\semi\fontsize{\fpeval{#1*26}}{\fpeval{#1*26}}\selectfont\addfontfeature{LetterSpace=-2.5}\color{#2}cotalia}}}

% ---------------------------------------------------------------- encabezado y pie
\pagestyle{fancy}
\fancyhf{}
\renewcommand{\headrulewidth}{0pt}
\fancyhead[L]{\raisebox{-1pt}{\CotaliaMarca{0.13}{tinta}{cobre}}\hspace{5pt}{\semi\fontsize{8.5}{10}\selectfont\color{tinta}\addfontfeature{LetterSpace=-1}cotalia}}
\fancyhead[R]{\eyebrow[pizarra]{Etapa 01 \,/\, Propuesta de servicios}}
\fancyfoot[L]{\eyebrow[pizarra]{Finca Académica EARTH \,/\, Lote de cultivo \AreaHa~ha}}
\fancyfoot[R]{{\monomed\fontsize{7.5}{9}\selectfont\color{tinta}\numpag}{\monoreg\fontsize{7.5}{9}\selectfont\color{pizarra}\,/\,\pageref{LastPage}}}
\fancypagestyle{plain}{}

% ---------------------------------------------------------------- tablas, listas, leyendas
\newcommand{\cab}[1]{\leavevmode{\monomed\fontsize{6.6}{8}\selectfont\addfontfeature{LetterSpace=10}\color{pizarra}#1}}
\newcolumntype{L}[1]{>{\raggedright\arraybackslash}p{#1}}
\newcolumntype{C}[1]{>{\centering\arraybackslash}p{#1}}
\newcolumntype{Y}{>{\raggedright\arraybackslash}X}
\renewcommand{\arraystretch}{1.5}
\AtBeginEnvironment{table}{\fontsize{8.9}{12.4}\linespread{1}\selectfont}
\setlength{\heavyrulewidth}{0.7pt}
\setlength{\lightrulewidth}{0.3pt}
\setlength{\cmidrulewidth}{0.3pt}
\setlength{\aboverulesep}{0pt}
\setlength{\belowrulesep}{0pt}
\newcommand{\filete}{\arrayrulecolor{linea}\midrule\arrayrulecolor{tinta}}

\setlist[itemize]{label={\color{cobre}\rule[0.28em]{0.55em}{0.7pt}}, leftmargin=1.5em, itemsep=0.25em, topsep=0.3em}
\setlist[enumerate]{label={\monomed\small\color{cobre}0\arabic*}, leftmargin=2.2em, itemsep=0.35em, topsep=0.3em}

\DeclareCaptionFont{capetq}{\monomed\fontsize{6.8}{9}\selectfont\addfontfeature{LetterSpace=12}\color{cobre}}
\DeclareCaptionFont{captxt}{\fontsize{8.3}{12.2}\linespread{1}\selectfont\color{pizarra}}
\DeclareCaptionLabelFormat{mayus}{\MakeUppercase{#1}~#2}
\captionsetup{labelfont=capetq, textfont=captxt, labelformat=mayus, labelsep=quad,
  justification=raggedright, singlelinecheck=false, skip=8pt}
\captionsetup[figure]{name=Fig.}
\captionsetup[table]{name=Tabla, position=top, skip=6pt}

% tarjeta de cifra clave
\newcommand{\cifra}[3]{% #1 valor, #2 unidad, #3 etiqueta
  \begin{minipage}[t]{\linewidth}
    {\color{tinta}\rule{\linewidth}{0.7pt}}\par\vspace{6pt}
    {\xlight\fontsize{30}{32}\linespread{1}\selectfont\color{tinta}#1}{\light\fontsize{12}{14}\selectfont\color{pizarra}\,#2}\par\vspace{4pt}
    \eyebrow[pizarra]{#3}
  \end{minipage}}

% celda de característica (rejilla de fortalezas)
\newcommand{\rasgo}[2]{%
  \begin{minipage}[t]{\linewidth}
    {\color{linea}\rule{\linewidth}{0.4pt}}\par\vspace{5pt}
    {\semi\color{tinta}#1}\par\vspace{2pt}
    {\fontsize{8.8}{13}\linespread{1}\selectfont #2}
  \end{minipage}}

% estado de un entregable: completado (círculo lleno), parcial (medio), pendiente (vacío)
\newcommand{\marcaC}{\tikz[baseline=-0.55ex]\fill[tinta] (0,0) circle (0.95mm);}
\newcommand{\marcaP}{\tikz[baseline=-0.55ex]{\fill[cobre] (0,-0.95mm) arc[start angle=-90, end angle=90, radius=0.95mm] -- cycle;
  \draw[cobre, line width=0.6pt] (0,0) circle (0.95mm);}}
\newcommand{\marcaN}{\tikz[baseline=-0.55ex]\draw[pizarra, line width=0.6pt] (0,0) circle (0.95mm);}
\newcommand{\estC}{\leavevmode\marcaC\hspace{4pt}{\semi\color{tinta}Completado}}
\newcommand{\estP}{\leavevmode\marcaP\hspace{4pt}{\semi\color{cobre}Parcial}}
\newcommand{\estN}{\leavevmode\marcaN\hspace{4pt}{\color{pizarra}Pendiente}}
\newcommand{\entregable}[2]{\leavevmode{\semi\color{tinta}#1}\par{\monoreg\fontsize{6.8}{9}\selectfont\color{pizarra}#2}}

% casilla de avance por entregable (tira de progreso)
\tikzset{celda/.style={minimum width=1.28cm, minimum height=0.95cm, inner sep=0pt, font=\monomed\fontsize{9}{10}\selectfont}}
\newcommand{\celda}[3]{% #1 x, #2 número, #3 estado C/P/N
  \ifx#3C\node[celda, fill=tinta, text=white] at (#1,0) {#2};\fi
  \ifx#3P\node[celda, draw=cobre, line width=0.6pt, text=cobre,
      path picture={\fill[cobre] (path picture bounding box.south west) rectangle ([yshift=0.16cm]path picture bounding box.south);}] at (#1,0) {#2};\fi
  \ifx#3N\node[celda, draw=linea, line width=0.6pt, text=pizarra] at (#1,0) {#2};\fi}

% globo en color para EGM2008 (modelo del geoide, sin logotipo propio)
\definecolor{globo}{HTML}{2F74B5}
\newcommand{\globoicono}{%
  \begin{tikzpicture}[baseline=0pt, x=1pt, y=1pt]
    \fill[globo] (12.5,12.5) circle (12.5);
    \begin{scope}
      \clip (12.5,12.5) circle (12.5);
      \foreach \rx in {4.5,9} {\draw[white, line width=0.55pt] (12.5,12.5) ellipse [x radius=\rx, y radius=12.5];}
      \draw[white, line width=0.55pt] (12.5,0) -- (12.5,25);
      \foreach \y in {5.5,12.5,19.5} {\draw[white, line width=0.55pt] (0,\y) -- (25,\y);}
    \end{scope}
    \draw[cobre, line width=1.4pt] (12.5,12.5) circle (12.5);
  \end{tikzpicture}}

% tarjeta de tecnología: logotipo, nombre, versión y función
\newcommand{\pila}[4]{%
  \begin{minipage}[t]{\linewidth}
    \setlength{\parskip}{0pt}%
    {\color{linea}\rule{\linewidth}{0.4pt}}\par\vspace{10pt}
    {\color{tinta}\fontsize{25}{25}\selectfont\raisebox{0pt}[26pt][2pt]{#1}}\par\vspace{8pt}
    {\semi\color{tinta}#2}\hspace{5pt}{\monoreg\fontsize{7}{9}\selectfont\color{pizarra}#3}\par\vspace{3pt}
    {\fontsize{8.3}{12}\linespread{1}\selectfont #4}
  \end{minipage}}

% subtítulo menor dentro de una subsección: no queda solo al pie de la página
\newcommand{\subtema}[1]{\par\needspace{5\baselineskip}\eyebrow[pizarra]{#1}\par\vspace{2pt}}

% nota lateral
\newcommand{\nota}[1]{\leavevmode{\fontsize{8.3}{12.4}\linespread{1}\selectfont\color{pizarra}#1\par}}

\graphicspath{{../outputs/}{../}}
